% ============================================================================
%  Streszczenie i Abstract
% ============================================================================
\clearpage
\begingroup
\setstretch{1.2}
\noindent{\fontsize{16}{20}\selectfont\bfseries Streszczenie\par}
\vspace{10pt}

Przedmiotem pracy jest projekt i implementacja trójwymiarowej gry zręcznościowej
w~silniku Unity 2022.3 z~wykorzystaniem języka C\#. Gracz steruje
kulą, nadając jej kierunek i~siłę gestem przeciągnięcia, zbiera monety rozmieszczone na
planszy, omija przeciwników oraz przeszkody ruchome, a~po zebraniu wszystkich monet
dociera do punktu końcowego poziomu. Gra obejmuje trzy poziomy, system punktów życia
i~leczenia, pola przyspieszające, odbijacze, obracające się belki, ruchome ściany, losowy
wiatr, minimapę oraz pomiar czasu z~zapisem najlepszych wyników. Część inżynierska pracy
dotyczy architektury aplikacji: kompozycji zależności z~użyciem biblioteki Zenject, podziału
kodu na 50~zestawów kompilacji (\emph{assembly definitions}), oddzielenia logiki rozgrywki
od komponentów silnika, interfejsu użytkownika w~podejściu MVVM oraz obsługi zdarzeń
za pomocą biblioteki reaktywnej R3 i~operacji asynchronicznych UniTask. Parametry rozgrywki
są wczytywane z~plików JSON, a~efekty cząsteczkowe korzystają z~puli obiektów. Aplikację
zbudowano na platformę Android z~zapleczem IL2CPP. W~pracy opisano także testy
funkcjonalne oraz zmiany wprowadzone na podstawie sesji testowych rozgrywki.

\vspace{18pt}
\noindent{\fontsize{16}{20}\selectfont\bfseries Abstract\par}
\vspace{10pt}

\begin{otherlanguage}{english}
This thesis presents the design and implementation of a~3D arcade
game built with the Unity 2022.3 engine and C\#. The player launches a~ball with
a~drag-and-release gesture that sets its direction and power, collects the coins placed on
the level, avoids enemies and moving obstacles, and reaches the finish point that appears
once every coin has been collected. The game contains three levels, health and healing,
boost pads, bumpers, spinning bars, moving walls, random wind, a~minimap and a~level timer
with stored best times. The engineering part concerns the application architecture:
dependency injection with the Zenject framework, code split into 50 assembly definitions,
game logic kept apart from engine components, a~user interface following the MVVM pattern,
and event handling built on the reactive library R3 together with UniTask for asynchronous
operations. Gameplay parameters are loaded from JSON files and particle effects are reused
through an object pool. The application was built for Android using the IL2CPP backend.
The thesis also describes functional testing and the changes made after play-testing
sessions.
\end{otherlanguage}
\endgroup
\clearpage
